\newpage
\titledquestion{Splitting Trees}

\textbf{Splitting Trees}

Recall binary search trees of integers:
\begin{codeblock}
  datatype tree = Empty | Node of tree * int * tree
\end{codeblock}
A \code{tree} is a \textbf{BST} if at every \code{Node (l, x, r)}, every value in
\code{l} is less than \code{x} and every value in \code{r} is greater. Assume all
values are distinct.

We want to \textbf{split} a BST by a pivot \code{p} into two BSTs: one holding
every value less than \code{p}, the other every value greater than \code{p}
(if \code{p} itself appears, it is discarded).

\spec
  {split}
  {\code{int * tree -> tree * tree}}
  {the given \code{tree} is a BST}
  {\code{split (p, T)} $\eeq$ \code{(T1, T2)}, where \code{T1} is a BST of every
  value in \code{T} less than \code{p}, and \code{T2} is a BST of every value
  greater than \code{p}.}

Here is one correct implementation. It flattens the tree, partitions the values,
and rebuilds:
\begin{codeblock}
  fun insert (x, Empty) = Node (Empty, x, Empty)
    | insert (x, Node (l, y, r)) =
        if x < y then Node (insert (x, l), y, r)
        else if x > y then Node (l, y, insert (x, r))
        else Node (l, y, r)

  fun fromList xs = List.foldl insert Empty xs

  fun inorder Empty = []
    | inorder (Node (l, x, r)) = inorder l @ (x :: inorder r)

  fun splitNaive (p, T) =
    let
      val xs = inorder T
      val less = List.filter (fn y => y < p) xs
      val more = List.filter (fn y => y > p) xs
    in
      (fromList less, fromList more)
    end
\end{codeblock}

\begin{parts}
  \part[4]\hfill

    Assume \code{T} is a balanced BST with \code{n} nodes. Give the \textbf{work}
    of \code{splitNaive (p, T)} as a big-$O$ in terms of \code{n}, and briefly
    justify it. You may use that \code{inorder} on a balanced tree is
    $O(n \log n)$, \code{filter} is $O(n)$, and a single \code{insert} into a
    balanced tree of size $k$ is $O(\log k)$.

    \begin{solutionorbox}[14em]
      \code{inorder} is $O(n \log n)$; the two \code{filter}s are $O(n)$; and
      \code{fromList} does up to \code{n} inserts, each $O(\log n)$ into a tree
      of size $\le n$, for $O(n \log n)$. The total is
      \[ O(n \log n). \]
    \end{solutionorbox}

  \newpage

  \part[12]\hfill

    We can do much better --- in $O(\text{height})$ work --- by working directly
    on the tree and never flattening it.

    Implement \code{split} with the type and specification above, recursively on
    the tree, \textbf{without} using \code{inorder}, \code{fromList},
    \code{filter}, or any list operations.

    \begin{constraint}
      Your solution must run in $O(\text{height of } T)$ work. This means that
      at each node you visit, you may make at most \textbf{one} recursive call.
    \end{constraint}

    \textbf{Hint:} At \code{Node (l, x, r)}, compare \code{x} to \code{p}. What
    does the BST invariant already tell you about where \textit{all} of \code{l}
    belongs, or all of \code{r}? One of the two subtrees can be kept whole,
    untouched; only the other needs to be split.

    \begin{solutionorbox}[30em]
      \begin{codeblock}
        fun split (p, Empty) = (Empty, Empty)
          | split (p, Node (l, x, r)) =
              if x < p then
                let
                  val (rl, rr) = split (p, r)
                in
                  (Node (l, x, rl), rr)
                end
              else if x > p then
                let
                  val (ll, lr) = split (p, l)
                in
                  (ll, Node (lr, x, r))
                end
              else
                (l, r)
      \end{codeblock}

      When \code{x < p}, the invariant guarantees \code{l} and \code{x} are all
      below \code{p}, so they go straight into \code{T1}; only \code{r} can hold
      values on either side of \code{p}, so it is the one subtree we split. Its
      lesser part \code{rl} joins \code{l} and \code{x}; its greater part
      \code{rr} is \code{T2}. The \code{x > p} case is symmetric, and
      \code{x = p} discards \code{x} with \code{l}, \code{r} already the two
      halves. Only one recursive call is made per node, down a single path.
    \end{solutionorbox}

  \part[3]\hfill

    Give the \textbf{work} of your \code{split} on a balanced BST of \code{n}
    nodes, as a big-$O$ in terms of \code{n}, and state in one sentence how it
    compares to \code{splitNaive}.

    \begin{solutionorbox}[10em]
      Each recursive call descends one level and does $O(1)$ work, and there is
      one call per level down a single root-to-leaf path, so the work is
      $O(\log n)$ on a balanced tree.

      This is an exponential improvement over \code{splitNaive}'s $O(n \log n)$:
      by using the ordering invariant to keep whole subtrees untouched, we avoid
      ever looking at most of the tree.
    \end{solutionorbox}
\end{parts}
